\cotefile{89}
\scanpage{12}
Ceci étant, on \unsure{a} des équivalences de catégories

\diagram{\begin{tikzcd}[ampersand replacement=\&,column sep=small,row sep=large]
\& \bigl(\underline{\mathrm{Tors}}(V)\bigr)
  \arrow[dl,leftrightarrow,"\approx"'] \arrow[dr,leftrightarrow,"\approx"] \& \\
{\llap{\text{(7)}\qquad}\displaystyle\prod_{i\in\tau}\bigl(\mathrm{Tors}(D_i)\bigr)}
 \& \& {\displaystyle\prod_{i\in\tau'}\bigl(\mathrm{Tors}(D_i)\bigr).}
\end{tikzcd}}

Notons que pour une droite vectorielle $D$ sur $\mathbb F_3$, la catégorie $\mathrm{Tors}(D)$ est isomorphe canoniquement \unsure{à|:} celle des ens. $E$ de cardinal $3$, munis d'une bijection
\[
\underset{\substack{\text{(ens des 2}\\\text{orientations de }E,\\
\text{i.e. des deux}\\\text{permutations circulaires}\\
\text{de }E\text{ inverses}\\\text{l'une de l'autre}}}{\omega_E}
 \simeq D^{*}
\]
de sorte que la catégorie $\prod_{i\in\tau}\bigl(\mathrm{Tors}(D_i)\bigr)$ peut être regardée comme celle des familles à deux éléments, indexées par $\tau$, d'ensembles de cardinal $3$
\[
(\deleted{F}P_i)_{i\in\tau}
\]
munis chacun d'une bijection
\[
\omega_{P_i}\simeq D_i^{*}\qquad(i\in\tau)
\]
-- et de même pour la catégorie $\prod_{i\in\tau'}\bigl(\mathrm{Tors}(D_i)\bigr)$.
% archive sheet number: 11
